\section{The Frozen \nova{}-S Science R1 Release}
\label{app:nova}

\subsection{Provenance, immutability, and input identities}

NOVA-S Science R1 is the immutable internal V50/QE/FULL lineage applied to
the V45 WASP-17 benchmark cube. A release manifest authenticates 419 files
covering code, configuration, predictions, and diagnostics, and internal
version names are retained in the code so the certified run stays
byte-traceable; diagnostics and candidate models live outside the frozen
directory. The wavelength and trace solution is the current PASTASOSS
reference \citep{BainesEtAl2023Trace,BainesEtAl2023Wavelength} with the
WebbKernel implementation of the ATOCA response mapping
\citep{DarveauBernierEtAl2022,JWSTPipeline2025}; legacy SPECTRACE products
are not used. The runtime manifest pins the exact software releases,
including the JWST calibration pipeline \citep{JWSTPipeline2025} and the
jaxoplanet transit engine \citep{Jaxoplanet2025}.

Two calibration masks with different counts coexist by construction and
must not be conflated: the source-curvature calibration uses all 102
event-complement integrations, while the OOT baseline mask of the q-star build
retains 101 integrations because its timing and quality manifest excludes
one further sample; the robust template fit of that build then spans all
transit phases in stratified folds.

\subsection{Detector support and reporting bins}

The initial mapping contained 2{,}726 native (order, detector-column)
groups. Six edge groups without reliable mapping are excluded
(order-1 columns 23, 24, and 2043; order-2 columns 974, 1677, and 1678),
leaving 2{,}720 fitted groups; the q-star product records eight further
unsupported response groups handled by the runtime support logic.

PASTASOSS order-1 training anchors end at detector coordinate
$x=73.62893$, corresponding to $2.75977449\,\mu\mathrm{m}$ for this visit,
so detector columns 38--73 rely on polynomial extrapolation. The reporting
rule of Section~\ref{sec:methods} excludes every bin that touches these
columns. Bin 147 is 71.7\% inside the direct domain but is excluded as a
whole, and bin 148 lies outside it; bin 149 was already structurally
excluded from the release's 143 bins, so the rule removes two retained
bins and leaves 141. Bin 146 is the last retained bin and ends at
$2.740018\,\mu\mathrm{m}$.

\subsection{Positive shared-spectrum parameterisation}

In Equation~\ref{eq:depth-basis} the basis-width weights $w_k$ define the
weighted-zero-mean gauge of the morphology and $\bm{H}_v$ spans the
subspace orthogonal to the constant. The achromatic depth
$D_{\mathrm{global}}$ is bounded below by a floor of $10^{-6}$ ppm plus a
bound margin of $10^{-8}$ ppm in the internal transformed
parameterisation. The
optional two-coordinate order-discrepancy block is bounded at
$\pm 10^{-12}$: its terminal values sit at this numerical scale and carry
no physical freedom.

The orbital geometry is: period $3.73548546$~d;
$a/R_\star=17.59002069137651$; impact parameter $0.33859071787462686$;
mid-transit time $60023.52874613816$~$\mathrm{BJD_{TDB}}$; zero
eccentricity. The transit engine is the jaxoplanet Keplerian
implementation \citep{Jaxoplanet2025} with backend quadrature order 10 and
21-point finite-exposure quadrature. The limb-darkening priors use theory
scale 0.5 in logit space and curvature scale 2.0 per $\mu\mathrm{m}^2$
around the STAGGER reference \citep{MagicEtAl2015}.

\subsection{Q-star construction and known limitation}

Q-star is a normalised cross-dispersion profile within each native group;
it is not an additional spectrum. Its construction uses only the public
V45 cube. A per-pixel weighted OOT intercept-plus-slope baseline is
estimated; pixels qualify when at least 95\% of their OOT samples are
valid, their baseline is positive, and they lie above the 35th percentile
of eligible baseline brightness. A robust median of baseline-normalised
bright-pixel residuals defines a depth-free event template $\phi_o(t)$,
normalised after removal of constant and linear terms. Each pixel is then
fitted robustly as
\begin{equation}
 y_p(t)=a_{0,p}+a_{1,p}\,\tau(t)+A_p\,\phi_o(t)
\end{equation}
with Huber threshold 1.345 and four solve and reweight cycles. Per-group
amplitudes are normalised to unit sum, leaving total group brightness to
the fitted continuum. Folds alternate integrations within pre-OOT,
ingress, full-transit, egress, and post-OOT phases, so no fold sees only
one part of the event. A robust second-difference penalty smooths the
profile along detector column and integer trace-relative rank, with
strength chosen by opposite-fold chi-squared, followed by fixed A/B
stability gates: p95 half-L1 at most 0.15 and centroid difference at most
0.5 pixel.

The frozen profile carries one documented defect: the smoother works on
the integer trace-relative rank rather than the continuous trace-relative
coordinate. Audits of the release record 14 centroid resets larger than
0.5 pixel, a median 0.234-pixel registration difference relative to the
delivered source, and roughly 1--2\% excess in-mask red response amplitude
at the faintest columns. Because the continuum fixes each group's
brightness while q-star fixes its spatial profile, an over-amplitude of
this kind acts like a dilution and makes the recovered transit shallower
at the affected columns. The defect is documented rather than silently
repaired; its causal quantification belongs to the Results, and the rules
that any prospective repair must satisfy are given in
Appendix~\ref{app:benchmarking}.

\subsection{Detector uncertainties and the QE guard}

FITS ERR values are first multiplied by the frozen per-order scale $N_o$
(order 1: 3.0746905465; order 2: 4.1917779190). The QE arm then derives
one further scale for each detector row using only OOT data. Writing
$r(p)$ for the row of pixel $p$, alternating OOT folds fit a weighted
level plus slope and predict the opposite fold in both directions, and
\begin{equation}
 s_{r(p)}=\max\!\left[1,\;
 \sqrt{\frac{\sum r_{\mathrm{cv}}^2}{\max(n_{\mathrm{cv}}-4,\,1)}}\right],
\end{equation}
so rows whose OOT residuals exceed their nominal uncertainties are
down-weighted and no row is up-weighted. The final uncertainty is
$\sigma_{tp}=\mathrm{ERR}_{tp}\,N_o\,s_{r(p)}$. No wavelength information,
injected truth, or recovered spectrum enters this rule.

\subsection{E13b background construction}

Eight raw spatial maps comprise a standardised CRDS background anchor plus
seven broad polynomial surfaces ($1, y, x, xy, y^2, x^2, xy^2$),
orthonormalised on safe off-trace pixels. Candidate families and ranks 2,
4, and 8 were assessed by blocked off-trace cross-validation; rank 8 is
the smallest rank whose worst-fold RMSE lies within 5\% of the family
minimum. The eight spatial coefficients measured across OOT integrations
are centred and decomposed by singular value decomposition, and the first
of the leading eight temporal vectors with absolute transit-reference
correlation at most 0.25 is selected; the chosen component was index 1
with correlation $-0.0735$. The constant and this component give the two
temporal modes, whose product with the eight spatial modes yields the 16
additive coefficients. Their OOT expectation
$\bm{\alpha}_{\mathrm{off}}$ and normal matrix supply the quadratic anchor
of Equation~\ref{eq:varpro}. Recovery never subtracts a fixed background
and adds no further background variables.

\subsection{Common source-curvature selection}

On each 51-integration OOT fold, stable trace flux is integrated per order
after subtracting only the E13b OOT expectation for this calibration; a
weighted quadratic in time is divided by its weighted affine
approximation, and positive candidates are normalised to unit
training-fold mean and evaluated on the opposite fold. Held-out
standardised chi-squared was 3.26065 for the null model, 3.21178 for the
common factor, and 3.15235 for the per-order factor. The gates required
the common factor to improve the null by at least 1\%, and the per-order
factor additionally to improve the common one by at least 2\% and to
improve each order individually; the per-order candidate fell short of the aggregate requirement (a
1.85\% improvement over the common score) and its order-2 score (1.0226)
was worse than the common counterpart (0.99964), so the common factor was
selected. The frozen factor has unit OOT mean, range
0.999866573--1.000079673, event-mean correction $-111.390$ ppm, maximum
excursion 133.427 ppm, and fold correlation 0.999921.

\subsection{Exact profiled robust solve}

For fixed nonlinear state $\bm{\theta}$ and Huber weights, the linear
nuisance problem of Equation~\ref{eq:varpro} is solved exactly: the
equilibrated sparse continuum normal block is factorised and the
background coupling reduces to a 16-dimensional Schur complement. Write
the weighted augmented residual as
$\bm{r}(\bm{\theta})=\bm{C}(\bm{\theta})\widehat{\bm{\eta}}-\bm{d}$, where
$\bm{C}$ and $\bm{d}$ include the detector standardisation and the anchor
rows. Differentiating the normal equations for nonlinear coordinate
$\theta_a$, with $\bm{C}_a=\partial\bm{C}/\partial\theta_a$, gives the
movement of the inner optimum,
\begin{equation}
 (\bm{C}^{\mathsf T}\bm{C})\,\bm{\eta}_a
 =-\left(\bm{C}_a^{\mathsf T}\bm{r}
 +\bm{C}^{\mathsf T}\bm{C}_a\widehat{\bm{\eta}}\right),
\end{equation}
and the exact profiled derivative,
\begin{equation}
 \bm{r}_a=\bm{C}_a\widehat{\bm{\eta}}+\bm{C}\,\bm{\eta}_a .
\end{equation}
Both stellar and background coefficients respond to $\bm{\theta}$ through
these equations; treating them as fixed would omit part of the derivative
and can change the nonlinear step. Sufficient statistics are streamed, so
no dense detector Jacobian is formed.

Robustness uses the Huber loss \citep{Huber1964} on the standardised
residual $z$ with threshold $\delta=1.345$,
\begin{equation}
 \rho_{\delta}(z)=
 \begin{cases}
  \tfrac{1}{2}z^2, & |z|\leq\delta,\\
  \delta\left(|z|-\tfrac{1}{2}\delta\right), & |z|>\delta,
 \end{cases}
\end{equation}
implemented as exact outer iteratively reweighted least squares: the
weights are updated outside the inner least-squares solve, so the inner
problem stays linear and the combination is an exact M-estimator rather
than a mixture of incompatible robust losses. The outer nonlinear solve is
bounded trust-region reflective least squares
\citep{BranchColemanLi1999}.

Convergence requires maximum absolute weight change below $0.001$,
relative objective change below $10^{-8}$, shared-spectrum changes below
fixed RMS tolerances, a scale-aware Karush-Kuhn-Tucker measure below
$10^{-6}$, first-order optimality below $0.002$, agreement of the
deterministic starts, and a final confirmation fit at the recomputed
terminal weights. Soft and hard cycle caps are 25 and 50. After each
trust-region termination, a deterministic material-descent audit evaluates
the executable objective along scale-normalised Cauchy and Gauss-Newton
directions at linearised residual radii 0.03, 0.1, 0.3, and 1.0; a
physically valid improvement beyond the material threshold restarts the
same solver. The ordinary threshold and restart cap are $10^{-8}$ and 6,
and the final confirmation uses $10^{-10}$ and 12. An accepted fit must
show no tested material feasible descent.

\subsection{Deterministic starts and evaluation order}

Start 0 uses the white-light mean depth recorded in the release
configuration, zero chromatic morphology, and the theory limb-darkening
reference. Start 1 uses the same mean depth,
a predetermined weighted-zero-mean 250 ppm morphology perturbation, and
predetermined small limb-darkening perturbations. The starts are
independent and exchange no optimiser state; both converged at robust
cycle 14 in the frozen release.

Truth values are unavailable to all preprocessing, fitting, starting,
response selection, mask construction, and convergence code. Once a
terminal detector prediction is immutable, an evaluation sidecar computes
absolute RMSE, mean bias, and morphology RMSE. Because the fitted
objective is the detector-space robust likelihood, intermediate truth
error need not fall monotonically even while the fitted objective
improves.

\subsection{Frozen numerical constants}

\begin{table}[t]
 \centering
 \footnotesize
 \caption{Frozen constants of the NOVA-S Science R1 release.}
 \label{tab:frozen-constants}
 \setlength{\tabcolsep}{3pt}
 \begin{tabularx}{0.94\columnwidth}{@{}>{\raggedright\arraybackslash}p{0.46\columnwidth}>{\raggedright\arraybackslash}X@{}}
  \toprule
  Quantity & Frozen value \\
  \midrule
  Integrations (total / event / complement) & 229 / 127 / 102 \\
  Retained pixels per integration & 97{,}557 \\
  Native groups (initial / fitted) & 2{,}726 / 2{,}720 \\
  Shared spectral coefficients $K$ & 160 \\
  Continuum coefficients & 5{,}440 \\
  Background coefficients & $8\times2=16$ \\
  Limb-darkening coordinates & 8 \\
  ERR scale $N_o$, order 1 / order 2 & 3.0746905465 / 4.1917779190 \\
  Background temporal-correlation cap & 0.25 (selected $-0.0735$) \\
  Huber threshold $\delta$ & 1.345 \\
  IRLS cycle caps (soft / hard) & 25 / 50 \\
  Material-descent radii & 0.03, 0.1, 0.3, 1.0 \\
  Material threshold / cap (ordinary) & $10^{-8}$ / 6 \\
  Material threshold / cap (final) & $10^{-10}$ / 12 \\
  Exposure quadrature points & 21 \\
  Orbital period & $3.73548546$ d \\
  $a/R_\star$; impact parameter & 17.59002069; 0.33859072 \\
  Mid-transit time ($\mathrm{BJD_{TDB}}$) & 60023.52874614 \\
  Curvature event-mean correction & $-111.390$ ppm \\
  Reporting bins (release / paper rule) & 143 / 141 \\
  Last retained bin edge & $2.740018\,\mu\mathrm{m}$ \\
  \bottomrule
 \end{tabularx}
\end{table}

Table~\ref{tab:frozen-constants} collects the frozen constants referenced
from Section~\ref{sec:methods}; full-precision orbital values appear in
the text above, and the release manifest is the authoritative record.

\subsection{Rejected or unpromoted model classes}

The following classes were tested and are retained as negative results:
independent or priored per-order depth discrepancy; per-column priors on
continuum coefficients; aggregate order-contrast priors;
residual-location recentering with wavelength slopes; constant-recentering
families superseded by the frozen curvature factor; OOT-median source
profiles vulnerable to background contamination; extra background temporal
modes; a global narrow-response swap; local wavelength-specific smoothing
patches; recovery-side explicit order-overlap modelling; local
transit-amplitude templates at the red end; and continuous-coordinate
response candidates that improved red morphology while degrading order 2
or the headline gates. None of these is part of NOVA-S; each is recorded
with the evidence that rejected it, so failed ideas are neither repeated
nor quietly omitted.
